%% Shared abstract body, inputted by both main.tex (AAAI) and main_acl.tex (*ACL).
Large language model agents are increasingly deployed for financial information retrieval, and hence for downstream analytics and investment decisions. However, financial terminology is highly specific, and a question without sufficient specification may admit multiple defensible readings with materially different values, e.g., Meta Platforms' ``operating income'' is \$46.75B consolidated or \$62.87B for the Family of Apps segment, while each reading's ground truth is exact and programmatically verifiable. A capable agent should recognize the ambiguity and ask a clarification question instead of committing to a plausible but unintended reading. Existing financial benchmarks cannot see this, because they attach only one gold answer per question and therefore fail to distinguish agents that resolve the ambiguity from those making ungrounded guesses, a blind spot we call the \emph{single-gold illusion}. To address this illusion, we release \fininteract, a bilingual (English/Chinese) benchmark of 173 instances that pairs each question with a \emph{default} and an \emph{intended} interpretation (fixed by a disambiguating context) across a five-category ambiguity taxonomy, and grades whether an agent elicits and then integrates the right clarification. Re-grading the \emph{same} retrieval-augmented outputs against the intended rather than the default reading cuts GPT-4o's measured accuracy by a factor of $3.1$ ($37.0\%$ to $12.1\%$), so a single-gold convention would have reported roughly three times the model's true competence. Our fine-grained analysis shows that this failure is one of interaction: models answer above $90\%$ correctly once the interpretation is supplied, but reach at most $28.9\%$ under their own clarification at full scale ($44.0\%$ in a cross-vendor pilot), so the bottleneck lies in issuing the correct clarification and then integrating the answer. We further report a skill metric that localizes failure modes of different financial ambiguities across two well-powered categories (entity scope and metric definition) and three exploratory ones.
